% Nouveau EGA V (cote 25): archive cote 25, site page 137 (PDF page 138).
% typescript. Draft compared: none.
% AI-transcribed by Claude Opus 5.5 (claude-opus-5-5, adaptive thinking, effort high), Message Batches API, batch ega25b (msgbatch_01A7n5YuXubUPNK1zsepi8RG), 2026-11-13; not yet checked by eye.
\scanpage{25/137}
\mspage{\hand{40}}

IV 21 Faisceaux inversibles et diviseurs relatifs sur les fibrés projec-%
tifs et multiprojectifs; systêmes linéaires de diviseurs.

1. Faisceaux inversibles sur \xout{lesxf} un fibré projectif; application aux
morphismes d'un fibré projectif dans un autre.

2. Représentabilité de $\underline{\mathrm{Div}}^{L}_{X/S}$ ; diviseurs \inserted[typed]{relatifs} sur les fibrés projectifs et
multiprojectifs.

3. Systêmes linéaires de diviseurs$_{\mathrm{x}}$ et morphismes dans les fibrés projec-%
tifs.

4. Systêmes linéaires de diviseurs$_{\mathrm{x}}$ et Modules inversibles.

Les résultats développés dans le présent paragraphe et les suivants
sont déjà en partie de nature globale, et donnent notamment des complé-%
ments desur les schémas projectifs, utilisant \xout{de façon essentielle} les
constructions globales du Chap II, et même ici certains résultats du
Chap III, en plus des résultats purement locaux du présent Chapitre IV.

Un des buts du présent paragraphe est de développer le langage des
``systêmes linéaires de diviseurs'', lié d'une part à la classification des
morphismes dans un fibré projectif, d'autre part à celle des Modules inver
sibles sur un préschéma relatif donné. Notons d'ailleurs ici que les
``schémas de paramètres'' vraiment naturels pour les systêmes linéaires de
diviseurs sont les schémas de \uline{Brauer-Séveri}, qui généralisent les fibrés
projectifs, et peuvent se définir par exemple comme les fibrés qui devien-%
nent isomorphes à un fibré projectif après extension étale surjective de
la base. Comme leur étude utilise la théorie de la descente, développée
au Chap V, et que d'ailleurs leur classification équivaut à celle des \xout{fibr}
\xout{\unsure{fibrés principaux homogènes}} torseurs de groupe le groupe projectif, nous
renvoyons pour l'étude de ces schémas et ses liens avec la notion de sys-%
tême linéaire de diviseurs au Chap VI$_{\mathrm{x}}$ de notre savant ouvrage.

\mspage{\hand{137}}
